% Recuperación del producto nativo y su selector, no incorporación de RH.
\subsection{Multiplicative irreducibility and return periods}
\label{primos-retornos}

The distinction between structure, evaluation, and phase also appears in the
multiplicative sector. The local APP product contains both the residue
and the quotient; its iteration makes it possible to construct normal forms and
factorizations before examining their decimal periods. This development
places the expression \emph{prime mode} in a precise domain and relates
irreducibility to the operations already defined, not to an isolated
coincidence of periods (see \ref{primos-retornos}).

Let
\[
 \mathcal W_9=\{\varnothing\}\sqcup
 \bigsqcup_{\ell\geq1}\{1,\ldots,9\}^{\ell}.
\]
Sequences are ordered from the least significant digit. The evaluation
\(\operatorname{ev}_9(u)=\sum_j u_j9^j\), with
\(\operatorname{ev}_9(\varnothing)=0\), is a bijection onto the nonnegative
integers: for a positive integer, dividing \(n-1\) by \(9\)
uniquely determines the first digit and the continuation. This is
bijective nonadic numeration; it preserves the boundary representative
\(9\).

The sum \(\boxplus_\#\) is obtained by applying the additive APP sheet to the
digits at each position and propagating their quotients. If there is only one digit,
it constitutes the provisional residue; a nonzero incoming quotient is
normalized through a second additive operation. The sum of the
quotients is transported to the next position. To multiply a
sequence by a digit \(d\), the multiplicative cell produces
\(u_jd=9q_j+r_j\), and the incoming quotient \(c_j\) is incorporated through
\[
 (w_j,c_{j+1})=
 \begin{cases}
 (r_j,q_j),&c_j=0,\\
 \bigl(R^+(r_j,c_j),q_j+q^+(r_j,c_j)\bigr),&c_j>0.
 \end{cases}
\]
The procedure begins with \(c_0=0\) and incorporates the final quotient when it is
nonzero. During addition, \(c_j\leq2\); during multiplication by a digit,
\(c_j\leq9\). The local operations thus receive arguments within
their domain. Denote this second procedure by \(d\odot_\#u\),
with \(d\odot_\#\varnothing=\varnothing\).

For \(v=(v_0,\ldots,v_{m-1})\), the full product is constructed by
the recurrence
\[
 a_m=\varnothing,\qquad
 a_j=(9\odot_\#a_{j+1})\boxplus_\#(v_j\odot_\#u),
 \qquad u\boxtimes_\#v=a_0.
\]
The rule composes the two local sheets. Its evaluation satisfies
\begin{equation}
 \operatorname{ev}_9(u\boxtimes_\#v)
   =\operatorname{ev}_9(u)\operatorname{ev}_9(v).
 \label{eq:primos-entrelazamiento}
\end{equation}
Indeed, each transition preserves the partial sum plus the pending
quotient weighted by the power of \(9\). The recurrence then gives
\(\operatorname{ev}_9(a_j)=9\operatorname{ev}_9(a_{j+1})+
v_j\operatorname{ev}_9(u)\), whose iteration proves
\eqref{eq:primos-entrelazamiento}. Uniqueness of the normal form transports
associativity and commutativity to the constructed product; its unit is \((1)\).

\begin{definition}[Reduced multiplicative coproduct]
On the vector space of finite support generated by the nonempty
normal forms, define
\[
 \widetilde\Delta_\# e_w
 =\sum_{\substack{u\boxtimes_\#v=w\\u,v\ne(1)}}e_u\otimes e_v.
\]
A nonunit normal form is irreducible when its image under
\(\widetilde\Delta_\#\) is zero.
\end{definition}

The factorization fibres are finite by
\eqref{eq:primos-entrelazamiento}. In the completion
\(\mathcal H_\#=\ell^2(\mathcal W_9\setminus\{\varnothing\})\),
the images of distinct forms have disjoint supports. If
\(d_\#(w)\) counts the ordered nonunit factorizations, the domain
of the closure is
\[
 \operatorname{Dom}(\overline{\widetilde\Delta_\#})
 =\left\{x\in\mathcal H_\#:
       \sum_w d_\#(w)|x_w|^2<\infty\right\}.
\]
Indeed, the squared norm of the image is that weighted sum. Therefore,
the operator
\[
 A_\#=(\overline{\widetilde\Delta_\#})^*
                  \overline{\widetilde\Delta_\#}
\]
is diagonal, positive, and self-adjoint: its eigenvalue at \(e_w\) is the
number of ordered nonunit factorizations of \(w\). The projection
\[
 P_{\rm irr}=(I-|e_{(1)}\rangle\langle e_{(1)}|)
                    \mathbf1_{\{0\}}(A_\#)
\]
selects exactly the irreducible forms. Through multiplicative
evaluation, these correspond to the ordinary primes. Selection
depends on the factorization structure and keeps the unit
state separate.

For a value \(n\geq2\) coprime to \(10\), decimal division induces the
residue permutation \(r\mapsto10r\pmod n\). Return of the residue
\(1\) has period
\[
 \tau_n=\operatorname{ord}_n(10).
\]
If \(n=\prod_p p^{a_p}\), the Chinese remainder theorem gives
\[
 \tau_n=\operatorname{mcm}_{p^{a_p}\parallel n}
                    \operatorname{ord}_{p^{a_p}}(10).
\]
The composite return synchronizes the returns of the prime powers.
For a prime \(p\notin\{2,3,5\}\), we have
\[
 \tau_p\mid p-1,\qquad
 M_p=\frac{10^{\tau_p}-1}{p}\in\mathbb N,
 \qquad M_p\equiv0\pmod9.
\]
The last congruence expresses the nonadic closure of the decimal repetition:
\(10^{\tau_p}-1\) is a multiple of \(9\), and \(p\) is invertible modulo
\(9\). The same congruence holds for any \(n\) coprime to
\(90\). In particular, \(7,13,77,91,143\) have decimal period \(6\),
although only the first two are irreducible. The coproduct and the period
therefore record different properties of the same constructed value.

This distinction makes precise the relationship with aperiodic refinement.
The vacancy sequence determines the increments of the capacity
index; factorization determines the multiplicative modes;
decimal division determines their return times. The nonadic
representative preserves phase closure, while the quotient and
factorization preserve information that this closure does not distinguish. The
link with the spectral constructions of the treatise resides in those
operators and their intertwinings. The proofs of their subsequent
analytic results retain their own development; they are not replaced by
the periodicity of a residual observation.

\subsubsection{Vacancy discrepancy and Dirichlet readouts}

The link also admits an exact analytic expression. The same
encoding \(z_n=z_n(0)\), for \(n\geq1\), is preserved, and
two subsequent observables are considered: the series over all indices and the
product over the indices selected by irreducibility. For
\(\operatorname{Re}s>1\),
\[
 Z_{\rm vac}^{+}(s)=\sum_{n\geq1}\frac{z_n}{n^s},\qquad
 Z_{\rm vac}^{\times}(s)=\prod_p(1-p^{-s})^{-z_p}.
\]
Both converge absolutely in that half-plane. Their common factor is the
density \(\delta\), but their discrepancies differ.

\Needspace{21\baselineskip}
\begin{proposition}[Separation of vacancy density]
We have
\[
 Z_{\rm vac}^{+}(s)=\delta\zeta(s)+H_{\rm vac}(s),
 \qquad H_{\rm vac}\in\mathcal O(\operatorname{Re}s>0).
\]
In \(\operatorname{Re}s>1\), with logarithms defined by the absolutely
convergent Euler series,
\[
 Z_{\rm vac}^{\times}(s)=\zeta(s)^\delta G_{\rm vac}(s),\qquad
 \log G_{\rm vac}(s)=
 \sum_p(z_p-\delta)\log(1-p^{-s})^{-1}.
\]
\end{proposition}
\begin{proof}
The centred partial sums satisfy
\[
 A_N:=\sum_{n=1}^N(z_n-\delta)
 =\lfloor(N+1)\delta\rfloor-N\delta,\qquad |A_N|<1.
\]
Abel summation represents the centred term by
\(s\int_1^\infty A_{\lfloor x\rfloor}x^{-s-1}\,dx\). The integral
converges locally uniformly in \(\operatorname{Re}s>0\),
proving holomorphy. The second identity follows by separating
\(z_p=\delta+(z_p-\delta)\) in the logarithm of the product, within
its domain of absolute convergence.
\end{proof}

The decomposition makes visible what the readout on primes adds: it examines
the same vacancy sequence on a subset determined by
factorization. Control of the discrepancy over all integers gives
the first holomorphic extension; the second observable preserves a
discrepancy of its own on primes. These are two applications of the same
discrete structure, with explicit analytic domains. This articulation
explains their relevance to the arithmetic core without transferring to this
article the proofs of the terminal results of the treatise.
